\chapter{扭化}
\markboth{第 6 章\ 扭化}{第 6 章\ 扭化}

本节中，我们不假设所考虑的预叠与 DG 概形局部几乎有限型。我们将在 §6.7 恢复这个假设。

\section{Gerbe}

\ob{6.1.1}{n:6.1.1}记 $\mathrm{pt}/G_m$ 为群 $G_m$ 的分类叠。换言之，$\mathrm{pt}/G_m$ 是表示下述函子的代数叠：它把仿射 DG 概形 $S$ 对应到 $S$ 上线丛的 $\infty$-群胚。

事实上，由于 $G_m$ 是 Abel 群，叠 $\mathrm{pt}/G_m$ 有自然的 Abel 群结构。$\mathrm{pt}/G_m$ 上的乘法映射表示线丛的张量积。这个结构把 $\mathrm{pt}/G_m$ 升级为从仿射 DG 概形到 $\infty$-Picard 群胚（即连接谱）的函子。

对我们的目的而言，$G_m$-gerbe 是 $\mathrm{pt}/G_m$-torsor 的预层 $G$，它满足下述三个（互不等价的）条件之一：

(i) $G$ 在平展拓扑下局部非空\footnote{由 Toën 定理，这等价于在 fppf 拓扑下局部非空。}。

(ii) $G$ 在 Zariski 拓扑下局部非空。

(iii) $G$ 整体非空。

具体地说，记 $B^{\mathrm{naive}}(\mathrm{pt}/G_m)$ 为 $\mathrm{pt}/G_m$ 的分类预叠。它由单纯预叠
\[
B^{\mathrm{naive}}(\mathrm{pt}/G_m):=\cdots\mathrm{pt}/G_m\times\mathrm{pt}/G_m\rightrightarrows\mathrm{pt}/G_m\to\mathrm{pt}
\]
的几何实现给出。记 $B^{\mathrm{Zar}}(\mathrm{pt}/G_m)$（相应地，$B^{\mathrm{et}}(\mathrm{pt}/G_m)$）为预叠 $B^{\mathrm{naive}}(\mathrm{pt}/G_m)$ 的 Zariski（相应地，平展）层化。

预叠 $B^{\mathrm{et}}(\mathrm{pt}/G_m)$、$B^{\mathrm{Zar}}(\mathrm{pt}/G_m)$ 与 $B^{\mathrm{naive}}(\mathrm{pt}/G_m)$ 分别表示满足上述条件 (i)、(ii)、(iii) 的 $G_m$-gerbe。

记 $(\mathbf{Ge}_{G_m})_{\DGSch_{\mathrm{aff}}}$ 为函子
\[
(\DGSch_{\mathrm{aff}})^{\mathrm{op}}\to\infty\text{-}\mathrm{PicGrpd},
\]
它把仿射 DG 概形 $S$ 对应到 $G_m$-gerbe 的群胚，其中 gerbe 取上述三种概念中的任何一种。

\begin{othm}{注}{6.1.2}{n:6.1.2}
虽然这三个版本并不给出等价的 $G_m$-gerbe 概念，但我们很快会看到，它们导致相同的扭化定义，因为相关的 gerbe 将是那些限制到 $\clred S$ 上平凡化的 gerbe。
\end{othm}

\ob{6.1.3}{n:6.1.3}定义函子
\[
(\mathbf{Ge}_{G_m})_{\PreStk}:(\PreStk)^{\mathrm{op}}\to\infty\text{-}\mathrm{PicGrpd}
\]
为 $(\mathbf{Ge}_{G_m})_{\DGSch_{\mathrm{aff}}}$ 沿
\[
(\DGSch_{\mathrm{aff}})^{\mathrm{op}}\hookrightarrow(\PreStk)^{\mathrm{op}}
\]
的右 Kan 延拓。也就是说，对 $Y\in\PreStk$，
\[
\mathbf{Ge}_{G_m}(Y):=\varprojlim_{S\in(\DGSch_{\mathrm{aff}}/Y)^{\mathrm{op}}}\mathbf{Ge}_{G_m}(S).
\]
等价地，
\[
\mathbf{Ge}_{G_m}(Y)=\Maps_{\PreStk}\big(Y,B^{?}(\mathrm{pt}/G_m)\big)
\]
其中 $?=\mathrm{naive},\mathrm{Zar}$ 或 $\mathrm{et}$。

因此，粗略地说，$Y$ 上的 $G_m$-gerbe 是对每个映射到 $Y$ 的 $S\in\DGSch_{\mathrm{aff}}$ 指定 $S$ 上的一个 $G_m$-gerbe，并对 $S$ 函子性地相容。

对子范畴 $\mathcal{C}\subset\PreStk$，记 $(\mathbf{Ge}_{G_m})_{\mathcal{C}}$ 为限制 $(\mathbf{Ge}_{G_m})_{\PreStk}|_{\mathcal{C}}$。

\section{扭化的概念}

\ob{6.2.1}{n:6.2.1}设 $Y$ 是预叠。$Y$ 上扭化的 Picard 群胚定义为
\[
\Tw(Y):=\ker\big(p^*_{\dR,Y}:\mathbf{Ge}_{G_m}(Y_{\dR})\to\mathbf{Ge}_{G_m}(Y)\big),
\]
其中 $\mathbf{Ge}_{G_m}$ 按 naive、Zar 或 et 三种版本中的任何一种理解。我们很快会看到（见 §6.4），三个版本都等价。

粗略地说，$Y$ 上的扭化 $T$ 是如下数据：对每个带有映射 $\clred S\to Y$ 的 $S\in\DGSch_{\mathrm{aff}}$，指定一个对象 $G_S\in\mathbf{Ge}_{G_m}(S)$，它在映射 $S_1\to S_2$ 下相容地变化。此外，对上述映射到映射 $S\to Y$ 的每个延拓，指定 $G_S$ 的一个平凡化，它同样对映射 $S_1\to S_2$ 函子性地变化。

\begin{othm}{注}{6.2.2}{n:6.2.2}
当我们写 $\ker(A_1\to A_2)$（其中 $A_1\to A_2$ 是 $\infty$-PicGrpd 中的映射）时，指的是
\[
A_1\times_{A_2}\{*\},
\]
其中纤维积在 $\infty$-PicGrpd 中取。也就是说，这等于在所有（即不必连接的）谱的范畴中取的纤维积的连接截断。
\end{othm}

\ob{6.2.3}{n:6.2.3}例。设 $L$ 是 $Y$ 上的线丛。如下定义 $Y$ 上的扭化 $T(L)$：它对每个带有映射 $\clred S\to Y$ 的 $S\in\DGSch_{\mathrm{aff}}$ 指定平凡的 $G_m$-gerbe。对映射 $S\to Y$，把上述 gerbe 的平凡化乘以 $L$ 来平凡化它。

\ob{6.2.4}{n:6.2.4}显然扭化构成函子
\[
\Tw_{\PreStk}:\PreStk^{\mathrm{op}}\to\infty\text{-}\mathrm{PicGrpd}.
\]
对态射 $f:Y_1\to Y_2$，记 $f^*$ 为相应的函子
\[
\Tw(Y_2)\to\Tw(Y_1).
\]
若 $\mathcal{C}$ 是 $\PreStk$ 的子范畴（例如 $\mathcal{C}=\DGSch_{\mathrm{aff}}$ 或 $\DGSch$），记 $\Tw_{\mathcal{C}}$ 为 $\Tw_{\PreStk}$ 到 $\mathcal{C}^{\mathrm{op}}$ 的限制。

\ob{6.2.5}{n:6.2.5}按构造，函子 $\Tw_{\PreStk}$ 把 $\PreStk$ 中的余极限送到 $\infty$-PicGrpd 中的极限。于是由推论 1.1.5 得到：

\begin{othm}{引理}{6.2.6}{n:6.2.6}
函子 $\Tw_{\PreStk}$ 同构于 $\Tw_{\mathcal{C}}$ 沿
\[
\mathcal{C}^{\mathrm{op}}\hookrightarrow\PreStk^{\mathrm{op}}
\]
的右 Kan 延拓，其中 $\mathcal{C}$ 是下列范畴之一：
\[
\DGSch_{\mathrm{aff}},\qquad \DGSch_{\mathrm{qs\text{-}qc}},\qquad \DGSch.
\]
\end{othm}

具体地说，这个引理说明映射
\[
\Tw(Y)\to\varprojlim_{S\in(\DGSch_{\mathrm{aff}}/Y)^{\mathrm{op}}}\Tw(S)
\]
是同构（并且 $\DGSch_{\mathrm{aff}}$ 可以换成 $\DGSch_{\mathrm{qs\text{-}qc}}$ 或 $\DGSch$）。

粗略地说，这意味着：在预叠 $Y$ 上指定一个扭化，等价于对每个映射到 $Y$ 的仿射 DG 概形 $S$ 指定一族相容的扭化。

\section{变体：其他结构群}

\ob{6.3.1}{n:6.3.1}设 $S$ 是仿射 DG 概形。考虑 Picard 群胚
\[
\mathbf{Ge}^{/\mathrm{red}}_{G_m}(S):=\ker\big(\mathbf{Ge}_{G_m}(S)\to\mathbf{Ge}_{G_m}(\clred S)\big).
\]
记 $(\mathbf{Ge}^{/\mathrm{red}}_{G_m})_{\DGSch_{\mathrm{aff}}}$ 为所得的函子
\[
(\DGSch)^{\mathrm{op}}\to\infty\text{-}\mathrm{PicGrpd}.
\]

\ob{6.3.2}{n:6.3.2}按定义，可以把 $\mathbf{Ge}^{/\mathrm{red}}_{G_m}(S)$ 看作关于 Abel 群预层
\[
(\mathcal{O}^{\times})^{/\mathrm{red}}_S:\ker(\mathcal{O}^{\times}_S\to\mathcal{O}^{\times}_{\clred S})
\]
的 gerbe（取 §6.1.1 三种版本中的任何一种）。

\ob{6.3.3}{n:6.3.3}除了 $G_m$-gerbe，还可以考虑 $G_a$-gerbe。有函子
\[
(\mathbf{Ge}_{G_a})_{\DGSch_{\mathrm{aff}}}:(\DGSch_{\mathrm{aff}})^{\mathrm{op}}\to\infty\text{-}\mathrm{PicGrpd},
\]
它把仿射 DG 概形 $S$ 对应到 $S$ 上 $G_a$-gerbe 的群胚。

注意与 $G_m$-gerbe 的情形不同，§6.1.1 中讨论的三种 gerbe 概念对 $G_a$-gerbe 是等价的。这是由于对仿射 DG 概形 $S$，
\[
H^2_{\mathrm{Zar}}(S,G_a)=H^2_{\mathrm{et}}(S,G_a)=0.
\]
于是 $(\mathbf{Ge}_{G_a})_{\DGSch_{\mathrm{aff}}}$ 由 $B^2(G_a)$ 表示，后者是相应单纯预叠的几何实现。

\ob{6.3.4}{n:6.3.4}按定义，对仿射 DG 概形 $S$，有
\[
\mathbf{Ge}_{G_a}(S)=B^2\big(\Maps(S,G_a)\big)\simeq B^2\big(\Gamma(S,\mathcal{O}_S)\big).
\]
特别地，看作连接谱时，$\mathbf{Ge}_{G_a}(S)$ 有基域 $k$ 上模的自然结构。这把 $(\mathbf{Ge}_{G_a})_{\DGSch_{\mathrm{aff}}}$ 升级为函子
\[
(\DGSch)^{\mathrm{op}}\to\infty\text{-}\mathrm{PicGrpd}_k,
\]
其中 $\infty\text{-}\mathrm{PicGrpd}_k$ 表示连接谱中 $k$-模范畴。注意由 Dold--Kan 对应，
\[
\infty\text{-}\mathrm{PicGrpd}_k\simeq\Vect^{\leq 0}.
\]
定义函子
\[
(\mathbf{Ge}_{G_a})_{\PreStk}:\PreStk^{\mathrm{op}}\to\infty\text{-}\mathrm{PicGrpd}_k
\]
为函子 $(\mathbf{Ge}_{G_a})_{\DGSch_{\mathrm{aff}}}$ 沿
\[
(\DGSch_{\mathrm{aff}})^{\mathrm{op}}\hookrightarrow\PreStk^{\mathrm{op}}
\]
的右 Kan 延拓。

\ob{6.3.5}{n:6.3.5}如同 $G_m$-gerbe 的情形，可以考虑 Picard 群胚
\[
\mathbf{Ge}^{/\mathrm{red}}_{G_a}(S):=\ker\big(\mathbf{Ge}_{G_a}(S)\to\mathbf{Ge}_{G_a}(\clred S)\big),
\]
并记 $(\mathbf{Ge}^{/\mathrm{red}}_{G_a})_{\DGSch_{\mathrm{aff}}}$ 为所得的函子
\[
(\DGSch)^{\mathrm{op}}\to\infty\text{-}\mathrm{PicGrpd}_k.
\]
按定义，对仿射 DG 概形 $S$，$\mathbf{Ge}^{/\mathrm{red}}_{G_a}(S)$ 是关于连接谱预层
\[
\mathcal{O}^{/\mathrm{red}}_S:=\ker(\mathcal{O}_S\to\mathcal{O}_{\clred S})
\]
的 gerbe。具体地，
\[
\mathbf{Ge}^{/\mathrm{red}}_{G_a}(S)\simeq B^2\big(\Gamma(S,\mathcal{O}^{/\mathrm{red}}_S)\big).
\]

\ob{6.3.6}{n:6.3.6}回忆 [GR1, Sect.~6.8.8]：指数映射定义同构
\[
\exp:\mathcal{O}^{/\mathrm{red}}_S\to(\mathcal{O}^{\times})^{/\mathrm{red}}_S.
\]
于是得到：

\begin{othm}{推论}{6.3.7}{n:6.3.7}
指数映射定义函子的同构
\begin{equation}
\exp:(\mathbf{Ge}_{G_a})^{/\mathrm{red}}_{\DGSch_{\mathrm{aff}}}\to(\mathbf{Ge}_{G_m})^{/\mathrm{red}}_{\DGSch_{\mathrm{aff}}}\tag{6.1}
\end{equation}
对 $(\mathbf{Ge}_{G_m})^{/\mathrm{red}}_{\DGSch_{\mathrm{aff}}}$ 的三种版本（naive、Zar 或 et）中的任何一种都成立。
\end{othm}

于是，若把 $\mathbf{Ge}^{/\mathrm{red}}_{G_m}(S)$ 实现为平展或 Zariski 拓扑中的 gerbe，则这个范畴有平凡的 $\pi_0$ 与 $\pi_1$。换言之，仿射 DG 概形上任何这样的 gerbe 都整体非空，且任何自同构都非典范地同构于恒等。

\ob{6.3.8}{n:6.3.8}同构 (6.1) 赋予 $\mathbf{Ge}^{/\mathrm{red}}_{G_m}(S)$（看作连接谱）基域 $k$ 上模的结构。这把 $(\mathbf{Ge}^{/\mathrm{red}}_{G_m})_{\DGSch_{\mathrm{aff}}}$ 升级为函子
\[
(\DGSch)^{\mathrm{op}}\to\infty\text{-}\mathrm{PicGrpd}_k.
\]
定义函子
\[
(\mathbf{Ge}^{/\mathrm{red}}_{G_m})_{\PreStk}:\PreStk^{\mathrm{op}}\to\infty\text{-}\mathrm{PicGrpd}_k
\]
为函子 $(\mathbf{Ge}^{/\mathrm{red}}_{G_m})_{\DGSch_{\mathrm{aff}}}$ 沿
\[
(\DGSch_{\mathrm{aff}})^{\mathrm{op}}\hookrightarrow\PreStk^{\mathrm{op}}
\]
的右 Kan 延拓。

\ob{6.3.9}{n:6.3.9}按定义，对 $Y\in\PreStk$，
\[
\mathbf{Ge}^{/\mathrm{red}}_{G_m}(Y):=\varprojlim_{S\in(\DGSch_{\mathrm{aff}}/Y)^{\mathrm{op}}}\mathbf{Ge}^{/\mathrm{red}}_{G_m}(S).
\]
粗略地说，对 $Y\in\PreStk$，对象 $G\in\mathbf{Ge}^{/\mathrm{red}}_{G_m}(Y)$ 是对每个 $S\in\DGSch_{\mathrm{aff}}/Y$ 指定一个对象 $G_S\in\mathbf{Ge}^{/\mathrm{red}}_{G_m}(S)$，并对每个 $S'\to S$ 指定同构
\[
G_{S'}\simeq f^*(G_S).
\]
下述结果由定义推出：

\begin{othm}{引理}{6.3.10}{n:6.3.10}
对 $Y\in\PreStk$，自然映射
\[
\mathbf{Ge}^{/\mathrm{red}}_{G_m}(Y)\to\ker\big(\mathbf{Ge}_{G_m}(Y)\to\mathbf{Ge}_{G_m}(\clred Y)\big)
\]
是同构，其中
\[
\clred Y:=\mathrm{LKE}_{(\mathrm{red}\,\Sch_{\mathrm{aff}})^{\mathrm{op}}\hookrightarrow(\PreStk)^{\mathrm{op}}}(Y|_{\mathrm{red}\,\Sch_{\mathrm{aff}}}).
\]
\end{othm}

\section{扭化：重新表述}

我们将证明：扭化的概念可以用
\[
(\mathbf{Ge}^{/\mathrm{red}}_{G_m})_{\PreStk},\qquad(\mathbf{Ge}_{G_a})_{\PreStk},\qquad(\mathbf{Ge}^{/\mathrm{red}}_{G_a})_{\PreStk}
\]
代替 $(\mathbf{Ge}_{G_m})_{\PreStk}$ 来表述。

\ob{6.4.1}{n:6.4.1}考虑函子
\[
\Tw^{/\mathrm{red}},\ \Tw_a,\ \Tw^{/\mathrm{red}}_a:\PreStk^{\mathrm{op}}\to\infty\text{-}\mathrm{PicGrpd},
\]
它们分别由
\[
\Tw^{/\mathrm{red}}(Y):=\ker\big(p^*_{\dR,Y}:\mathbf{Ge}^{/\mathrm{red}}_{G_m}(Y_{\dR})\to\mathbf{Ge}^{/\mathrm{red}}_{G_m}(Y)\big),
\]
\[
\Tw_a(Y):=\ker\big(p^*_{\dR,Y}:\mathbf{Ge}_{G_a}(Y_{\dR})\to\mathbf{Ge}_{G_a}(Y)\big)
\]
以及
\[
\Tw^{/\mathrm{red}}_a(Y):=\ker\big(p^*_{\dR,Y}:\mathbf{Ge}^{/\mathrm{red}}_{G_a}(Y_{\dR})\to\mathbf{Ge}^{/\mathrm{red}}_{G_a}(Y)\big)
\]
给出。

有下述由指数映射与显然的遗忘函子给出的函子图：
\begin{equation}
\begin{tikzcd}[column sep=large, row sep=large]
 & \Tw^{/\mathrm{red}}_a \arrow[r, "\exp"] \arrow[d] & \Tw^{/\mathrm{red}} \arrow[d] \\
 & \Tw_a \arrow[r] & \Tw
\end{tikzcd}\tag{6.2}
\end{equation}

\begin{othm}{命题}{6.4.2}{n:6.4.2}
(6.2) 中的函子都是等价。
\end{othm}

\begin{proof}
由指数映射给出的函子是等价，见 §6.3.6。

来证明 (6.2) 中右竖直映射是等价。这其实是同义反复：

两个函子都是沿
\[
(\DGSch_{\mathrm{aff}})^{\mathrm{op}}\to(\PreStk)^{\mathrm{op}}
\]
的右 Kan 延拓，因此只需证明所讨论的映射在对象 $S\in\DGSch_{\mathrm{aff}}$ 上取值时为同构。

有：
\[
\Tw^{/\mathrm{red}}(S)=\mathbf{Ge}^{/\mathrm{red}}_{G_m}(S_{\dR})\times_{\mathbf{Ge}^{/\mathrm{red}}_{G_m}(S)}\{*\}\overset{\text{引理 6.3.10}}{\simeq}
\]
\[
\simeq\ker\big(\mathbf{Ge}_{G_m}(S_{\dR})\to\mathbf{Ge}_{G_m}(\clred(S_{\dR}))\big)\times_{\ker(\mathbf{Ge}_{G_m}(S)\to\mathbf{Ge}_{G_m}(\clred S))}\{*\}=
\]
\[
=\ker\big(\mathbf{Ge}_{G_m}(S_{\dR})\to\mathbf{Ge}_{G_m}(\clred S)\big)\times_{\ker(\mathbf{Ge}_{G_m}(S)\to\mathbf{Ge}_{G_m}(\clred S))}\{*\}\simeq
\]
\[
\simeq\mathbf{Ge}_{G_m}(S_{\dR})\times_{\mathbf{Ge}_{G_m}(S)}\{*\}=\Tw(S).
\]
(6.2) 中左竖直箭头是等价，证明类似。
\end{proof}

\ob{6.4.3}{n:6.4.3}作为命题 6.4.2 的推论，得到：

\begin{othm}{推论}{6.4.4}{n:6.4.4}
扭化的三种版本——naive、Zar 与 et——等价。
\end{othm}

此外：

\begin{othm}{推论}{6.4.5}{n:6.4.5}
函子 $\Tw:(\PreStk)^{\mathrm{op}}\to\infty\text{-}\Grpd$ 典范地升级为函子
\[
(\PreStk)^{\mathrm{op}}\to\infty\text{-}\mathrm{PicGrpd}_k.
\]
\end{othm}

\ob{6.4.6}{n:6.4.6}例。可以利用 $\Tw$ 上自然的 $k$-模结构来产生更多的扭化例子。设 $L$ 是 $Y$ 上的线丛，$T(L)$ 是 §6.2.3 中的扭化。现在，对 $a\in k$，$\Tw(Y)$ 上的 $k$-模结构给出新的扭化 $T(L^{\otimes a})$。

\begin{othm}{注}{6.4.7}{n:6.4.7}
注意：仿射 DG 概形 $X$ 上的扭化 $T$ 未必平凡，即使在 Zariski 或平展拓扑局部也未必。对任意带有映射 $S\to X_{\dR}$ 的 $S\in\DGSch_{\mathrm{aff}}$，相应的 $S$ 上 $G_m$-gerbe 可以非典范地平凡化；但这种平凡化不一定能对不同选择的 $S$ 相容地选取。

对光滑经典 $X$，这样的 gerbe 的例子可以由一个非平展局部正合的闭 $2$-形式给出（见 §6.5.4）。

然而注意，例 6.4.6 中描述的 gerbe 是 Zariski 局部平凡的，因为线丛具有相应的性质。
\end{othm}

\ob{6.4.8}{n:6.4.8}收敛性。现在断言：

\begin{othm}{命题}{6.4.9}{n:6.4.9}
函子 $\Tw:(\DGSch_{\mathrm{aff}})^{\mathrm{op}}\to\infty\text{-}\mathrm{PicGrpd}$ 收敛\footnote{收敛的概念见 §1.3.1。}。
\end{othm}

\begin{proof}
我们将证明函子 $\Tw^{/\mathrm{red}}$ 收敛。为此，只需证明函子
\[
S\mapsto\mathbf{Ge}^{/\mathrm{red}}_{G_m}(S_{\dR})\quad\text{与}\quad\mathbf{Ge}^{/\mathrm{red}}_{G_m}(S)
\]
收敛。

$\mathbf{Ge}^{/\mathrm{red}}_{G_m}((-)_{\dR})$ 的收敛性是显然的，因为这个函子只依赖于底层的约化经典概形。于是只需证明 $\mathbf{Ge}^{/\mathrm{red}}_{G_m}(-)$ 的收敛性。

有：
\[
\mathbf{Ge}^{/\mathrm{red}}_{G_m}(S)=\mathbf{Ge}_{G_m}(S)\times_{\mathbf{Ge}_{G_m}(\clred S)}\{*\}.
\]
因此只需证明函子 $\mathbf{Ge}_{G_m}(-)$ 收敛。后者由
\[
\mathbf{Ge}_{G_m}(-)=B^2\big(\Maps(S,G_m)\big)
\]
以及 $G_m$ 作为 DG 概形收敛推出。
\end{proof}

可以把命题 6.4.9 同义反复地重新表述为：

\begin{othm}{推论}{6.4.10}{n:6.4.10}
函子 $\Tw_{\PreStk}$ 同构于
\[
\Tw_{{}^{<\infty}\DGSch_{\mathrm{aff}}}:=\Tw_{\DGSch_{\mathrm{aff}}}|_{{}^{<\infty}\DGSch_{\mathrm{aff}}}
\]
沿
\[
({}^{<\infty}\DGSch_{\mathrm{aff}})^{\mathrm{op}}\hookrightarrow(\DGSch_{\mathrm{aff}})^{\mathrm{op}}\hookrightarrow(\PreStk)^{\mathrm{op}}
\]
的右 Kan 延拓。
\end{othm}

\begin{othm}{注}{6.4.11}{n:6.4.11}
可以利用命题 6.4.9 证明函子 $\mathbf{Ge}_{G_m}$（三种版本中任何一种）也收敛。
\end{othm}

\ob{6.4.12}{n:6.4.12}局部几乎有限型情形下的扭化。推论 6.4.10 蕴含：为知道 $Y\in\PreStk$ 上的扭化（当 $Y$ 局部几乎有限型时），我们「不需要知道」非局部几乎有限型的 DG 概形。

\begin{othm}{推论}{6.4.13}{n:6.4.13}
(a) 对 $Y\in\PreStk_{\mathrm{laft}}$，自然定义的映射
\[
\Tw(Y)\to\varprojlim_{S\in(({}^{<\infty}\DGSch_{\mathrm{aff}}^{\mathrm{aft}})/Y)^{\mathrm{op}}}\Tw(S)
\]
是等价。

(b) 函子 $\Tw_{\PreStk_{\mathrm{laft}}}$ 同构于 $\Tw_{{}^{<\infty}\DGSch_{\mathrm{aff}}^{\mathrm{aft}}}$ 沿包含
\[
({}^{<\infty}\DGSch_{\mathrm{aff}}^{\mathrm{aft}})^{\mathrm{op}}\hookrightarrow(\DGSch_{\mathrm{aff}}^{\mathrm{aft}})^{\mathrm{op}}\hookrightarrow(\PreStk_{\mathrm{laft}})^{\mathrm{op}}
\]
的右 Kan 延拓。
\end{othm}

\begin{proof}
把 $\Tw$ 换成任意收敛预叠 $(\DGSch_{\mathrm{aff}})^{\mathrm{op}}\to\infty\text{-}\Grpd$ 时，这都是对的。
\end{proof}

\begin{othm}{注}{6.4.14}{n:6.4.14}
由注 6.4.11 推出：函子 $\mathbf{Ge}_{G_m}$（三种版本中任何一种）看作预层时属于 $\PreStk_{\mathrm{laft}}$。事实上，在 naive 版本中这是显然的，因为 $\mathrm{pt}/G_m$ 属于 $\PreStk_{\mathrm{laft}}$。对 Zar 与 et 版本，这由 [GL:Stacks, Corollary 2.5.7] 推出：该结论说，在 $n$-连接的预叠语境中，「局部有限型」这一条件在层化后仍保持，只要我们限制到截断预叠。
\end{othm}

\section{扭化 Picard 群胚的确定}

可以利用用 $G_a$-gerbe 对扭化的描述，给出扭化群胚的上同调描述。

\ob{6.5.1}{n:6.5.1}de Rham 上同调。设 $Y$ 是预叠。回忆 $Y$ 的凝聚上同调定义为
\[
H(Y):=\Gamma(Y,\mathcal{O}_Y)=\Maps_{\QCoh(Y)}(\mathcal{O}_Y,\mathcal{O}_Y).
\]
定义 $Y$ 的 de Rham 上同调为 $Y_{\dR}$ 的凝聚上同调，即
\[
H_{\dR}(Y):=H(Y_{\dR})=\Maps_{\QCoh(Y_{\dR})}(\mathcal{O}_{Y_{\dR}},\mathcal{O}_{Y_{\dR}}).
\]
注意由于 $\QCoh(Y_{\dR})$ 是稳定 $\infty$-范畴，上述 $\Maps$ 给出一个（不必连接的）谱。

设 $X$ 是光滑经典概形。此时由 §5.5，有
\[
H_{\dR}(X)=\Maps_{\mathrm{D\text{-}mod}^l(X)}(\mathcal{O}_X,\mathcal{O}_X).
\]
特别地，我们对 de Rham 上同调的定义与光滑经典概形的通常定义一致。

\ob{6.5.2}{n:6.5.2}考虑表示 $G_a$-gerbe 的函子 $B^2(G_a)$。按定义，对预叠 $Y$，有连接谱的同构：
\[
\Maps(Y,B^2(G_a))\simeq\tau^{\leq 0}\Maps_{\QCoh(Y)}(\mathcal{O}_Y,\mathcal{O}_Y)[2]\simeq\tau^{\leq 0}(H(Y)[2]).
\]
于是由命题 6.4.2 得到：

\begin{othm}{推论}{6.5.3}{n:6.5.3}
对预叠 $Y$，扭化的群胚由
\[
\Tw(Y)\simeq\tau^{\leq 2}\Big(H_{\dR}(Y)\times_{H(Y)}\{*\}\Big)[2]
\]
给出。
\end{othm}

\ob{6.5.4}{n:6.5.4}现在设 $X$ 是光滑经典概形。此时有
\[
H_{\dR}(X)\simeq\Gamma(X,\Omega^\bullet),
\]
其中 $\Omega^\bullet$ 是 $X$ 上的 de Rham 微分复形。自然映射
\[
H_{\dR}(X)\to H(X)
\]
由投影映射 $\Omega^\bullet\to\mathcal{O}_X$ 的整体截面给出。因此有
\[
\Tw(X)\simeq\tau^{\leq 2}\Big(H_{\dR}(X)\times_{H(X)}\{*\}\Big)[2]\simeq\tau^{\leq 2}\Gamma\big(X,\Omega^\bullet\times_{\mathcal{O}_X}0\big)[2]\simeq
\]
\[
\simeq\tau^{\leq 2}\Gamma\big(X,\tau^{\leq 2}(\Omega^\bullet\times_{\mathcal{O}_X}0)\big)[2].
\]
复形 $\tau^{\leq 2}(\Omega^\bullet\times_{\mathcal{O}_X}0)$ 与复形
\[
\Omega^1\to\Omega^{2,\mathrm{cl}}
\]
等同，其中 $\Omega^{2,\mathrm{cl}}$ 是闭 $2$-形式的层（放在上同调次数 $2$），映射是 de Rham 微分。

于是得到：$X$ 上扭化的 Picard 群胚由
\[
\Tw(X)\simeq\tau^{\leq 2}\Gamma(X,\Omega^1\to\Omega^{2,\mathrm{cl}})[2]
\]
给出。特别地，对光滑经典概形，我们的扭化定义与 [BB] 中 TDO 的概念一致。

\begin{annote}
这与 [BB] 的经典结论一致：光滑经典 $X$ 上扭化的 Picard 群胚是 $\tau^{\leq 2}\Gamma(X,\Omega^1\to\Omega^{2,\mathrm{cl}})[2]$，即扭化的对象由闭 $2$-形式给出、同构由 $1$-形式给出；相应的扭 D-模范畴就是 TDO 上的模范畴（§7.5.1）。\hfill【来源：原书 §0.4.2、§6.5.4；[BB]（A）】
\end{annote}


\section{扭化与无穷小群胚}

\ob{6.6.1}{n:6.6.1}设
\[
Y_1\rightrightarrows Y_0
\]
是 $\PreStk$ 中的群胚对象，$Y_\bullet$ 为相应的单纯对象。回忆这个群胚对象被 $G_m$ 中心扩张的概念。（这里 $G_m$ 可以换成任意交换群对象 $H\in\PreStk$。）

按定义，$Y_1\rightrightarrows Y_0$ 被 $G_m$ 的中心扩张是 $\mathbf{Ge}_{G_m}(|Y_\bullet|)$ 的一个对象，连同一个其沿
\[
Y_0\to|Y_\bullet|
\]
的限制的平凡化。

\ob{6.6.2}{n:6.6.2}粗略地说，这样的中心扩张的数据是：$Y_1$ 上的线丛 $L$，其沿退化映射 $Y_0\to Y_1$ 的拉回被平凡化，并且对三个映射
\[
p_{1,2},\ p_{2,3},\ p_{1,3}:Y_2\to Y_1,
\]
给定同构
\[
p_{1,2}^*(L)\otimes p_{2,3}^*(L)\simeq p_{1,3}^*(L),
\]
并满足进一步的融贯条件。

\ob{6.6.3}{n:6.6.3}对 $Y\in\PreStk$，考虑其无穷小群胚
\[
Y\times_{Y_{\dR}}Y\rightrightarrows Y.
\]
按定义，$Y$ 上的扭化给出其无穷小群胚被 $G_m$ 的中心扩张。

反之，由引理 1.2.4 得到：

\begin{othm}{推论}{6.6.4}{n:6.6.4}
假设 $Y$ 经典形式光滑。则上述函子
\[
\Tw(Y)\to\{Y\text{ 的无穷小群胚被 }G_m\text{ 的中心扩张}\}
\]
是等价。
\end{othm}

\section{ind-概形上的扭化}

\ob{6.7.1}{n:6.7.1}设 $X$ 是 $\DGIndSch_{\mathrm{laft}}$ 的对象。我们将证明推论 6.6.4 的断言对 $X$ 成立：

\begin{othm}{命题}{6.7.2}{n:6.7.2}
函子
\[
\Tw(X)\to\{X\text{ 的无穷小群胚被 }G_m\text{ 的中心扩张}\}
\]
是等价。
\end{othm}

本小节余下部分专门证明命题 6.7.2。

\ob{6.7.3}{n:6.7.3}第 1 步。由推论 6.4.13(b)，需要证明下述事情：

对每个 $S\in(\DGSch_{\mathrm{aff}}^{\mathrm{aft}})/X_{\dR}$，在 $S$ 上给定一个 $G_m$-gerbe 并给定其拉回到 $S\times_{X_{\dR}}X$ 的平凡化，等价于在单纯预叠
\[
S^\bullet:=S\times_{X_{\dR}}(X^\bullet/X_{\dR})
\]
上给定一个 $G_m$-gerbe 并给定其在 $0$-单形上（即 $S\times_{X_{\dR}}X$ 上）的平凡化。

\ob{6.7.4}{n:6.7.4}第 2 步。注意单纯预叠 $\clred(S^\bullet)$ 是常值的，值为 $\clred S$。于是由引理 6.3.7，可以用 $G_a$-gerbe 代替 $G_m$-gerbe。

因此只需证明映射
\[
\Maps_{\QCoh(S)}(\mathcal{O}_S,\mathcal{O}_S)\to\Tot\Big(\Maps_{\QCoh(S^\bullet)}(\mathcal{O}_{S^\bullet},\mathcal{O}_{S^\bullet})\Big)
\]
在 $\Vect$ 中是同构。

\ob{6.7.5}{n:6.7.5}第 3 步。注意对任意 $X'\in\DGIndSch_{\mathrm{laft}}$，典范映射
\[
\Maps_{\QCoh(X')}(\mathcal{O}_{X'},\mathcal{O}_{X'})\to\Maps_{\IndCoh(X')}(\omega_{X'},\omega_{X'})
\]
是同构。这例如由 DG 概形的相应断言，即引理 5.2.5 推出。

因此只需证明映射
\[
\Maps_{\IndCoh(S)}(\omega_S,\omega_S)\to\Tot\Big(\Maps_{\IndCoh(S^\bullet)}(\omega_{S^\bullet},\omega_{S^\bullet})\Big)
\]
是同构。

\ob{6.7.6}{n:6.7.6}第 4 步。注意 $S^\bullet$ 与映射
\begin{equation}
S\times_{X_{\dR}}X\to S\tag{6.3}
\end{equation}
的 Čech 神经等同。

如同命题 3.1.3 的证明，所有 $S^i$ 都属于 $\DGIndSch$，且态射 (6.3) 是 ind-固有的且满。

于是所求断言由 $\IndCoh$ 在 DG ind-概形的 ind-固有满映射下的下降推出，见 [GR1, Lemma 2.10.3]。
